\bookchapter{บทที่ 21: เจาะลึกความปลอดภัยเครือข่ายมือถือกับความเป็นส่วนตัว}

\chapterintro{ทุกวันเราใช้มือถือเข้าเน็ต ชำระเงิน แชท รู้สึกสะดวกสบายมาก แต่คุณเคยคิดไหมว่า ประวัติแชท รหัสชำระเงิน รูปภาพ… ข้อมูลพวกนี้โป้งเปล่าอยู่บนเน็ตบ่อยแค่ไหน? พูดตรงๆ คนส่วนใหญ่ไม่เคยคิดเรื่องนี้เลย แต่พอเกิดเรื่องขึ้น — เช่นบัญชีโดนขโมย บัตรโดนรูด — สายเกินแก้ บทนี้มาคุยกันให้เข้าใจว่าจะทำให้ชีวิตดิจิทัลของคุณโปร่งใสน้อยลงยังไง}

\booksection{หนึ่ง VPN: ผ้าคลุมล่องหนในเครือข่าย}

VPN ชื่อเต็มคือเครือข่ายส่วนตัวเสมือน (Virtual Private Network) ลองนึกว่ามันสร้างอุโมงค์ลับเฉพาะที่เข้ารหัส ระหว่างบ้านคุณกับเว็บไซต์ที่คุณจะไป คุณเดินอยู่ในอุโมงค์ คนข้างนอกมองไม่เห็นว่าคุณถืออะไรจะไปทางไหน

\begin{itemize}
  \item \textbf{เข้าเน็ตปกติ}: มือถือคุณ เครือข่ายสาธารณะ เว็บเป้าหมาย เราเตอร์ สวิตช์ ที่ผ่านตรงกลาง วงจรไหนก็มีคนแอบมองข้อมูลที่คุณส่งได้ โดยเฉพาะบนไวไฟสาธารณะอย่างร้านกาแฟสนามบิน พูดหยาบๆ ก็คือโป้งเปล่านั่นแหละ
  \item \textbf{ใช้ VPN}: มือถือคุณ อุโมงค์เข้ารหัส VPN เซิร์ฟเวอร์ VPN เว็บเป้าหมาย ข้อมูลถูกห่อหุ้มด้วยการเข้ารหัส คนแอบมองเห็นแค่ตัวอักษรระเกะระกะ นอกจากนี้เว็บเป้าหมายที่เห็นคือ IP ของเซิร์ฟเวอร์ VPN ไม่ใช่ IP จริงของคุณ
\end{itemize}

แล้ว VPN มันทำอะไรได้จริง? พูดง่ายๆ: หนึ่ง\textbf{คุ้มครองความเป็นส่วนตัวของคุณบนไวไฟสาธารณะ}; สอง,\textbf{ข้ามข้อจำกัดพื้นที่ภูมิภาค}(เช่นดูเนื้อหาที่เปิดเฉพาะบางประเทศ); สาม,\textbf{กันผู้ให้บริการเครือข่าย (ISP) บันทึกนิสัยการเข้าเว็บของคุณ}แต่ฉันต้องเตือนสักประโยค:\textbf{อย่าใช้ VPN ฟรี!}ของฟรีข้างหลังมักแปลว่าข้อมูลของคุณก็คือสินค้าของเขา ฉันมีเพื่อนคนหนึ่งใช้ VPN ฟรีตัวหนึ่ง ผลมือถือโผล่โฆษณาขึ้นมางงๆ ภายหลังค้นพบว่า VPN ตัวนั้นแพ็กข้อมูลการเข้าเว็บของเขาขายให้ผู้ลงโฆษณาไปแล้ว ดังนั้นเลือก VPN ไม่ก็เลือกแบรนด์ใหญ่แบบเสียเงิน (เช่น NordVPN, ExpressVPN) ไม่ก็อย่าใช้เสียเลย

อีกอย่าง VPN ไม่ใช่สารพัดประโยชน์ มันคุ้มครองได้แค่ข้อมูลระหว่างส่ง ถ้าเว็บเป้าหมายเองไม่ปลอดภัย หรือมือถือคุณโดนมัลแวร์ไปแล้ว VPN ก็ช่วยไม่ได้

\booksection{สอง กับดักหวานของไวไฟสาธารณะ — อันตรายกว่าที่คุณคิด}

คุณอาจคิดว่า แค่ต่อ Wi-Fi ฟรี มันจะเสี่ยงแค่ไหน? เล่าเรื่องจริงให้ฟัง: นักวิจัยด้านความปลอดภัยเคยทำการทดลอง วางฮอตสปอต Wi-Fi ชื่อ Starbucks\_Free ไว้ที่สนามบิน ไม่ถึงชั่วโมงมีคนต่อกันเกินร้อยคน ที่จริงฮอตสปอตนี้นักวิจัยสร้างเอง\textbf{ดาต้าที่ไม่ได้เข้ารหัส}ที่ผ่านฮอตสปอตนี้ — เช่นเนื้อหาเว็บ HTTP ที่คุณเปิด การสื่อสารของแอปที่ไม่ได้เปิด TLS — พวกเขาดักไปดูได้จริงๆ

แต่ต้องจำเงื่อนไขสำคัญ: วันนี้ WeChat ธนาคารมือถือ Alipay ที่คุณใช้ประจำ ส่วนใหญ่ผ่านการเข้ารหัส HTTPS หมดแล้ว พูดให้เห็นภาพ แม้ผู้โจมตีดักแพ็กเกจข้อมูลได้ สิ่งที่เห็นก็เป็นตัวอักษรระเกะระกะ\textbf{อ่านเนื้อหาแชทหรือรหัสชำระเงินของคุณโดยตรงไม่ได้}แต่พวกเขายังรู้ว่าคุณเข้าเว็บไหนบ้าง (โดเมน) ถ้าแอปเก่าๆ หรือหน้าเว็บบางตัวยังส่งข้อมูลเป็น HTTP เปล่า เนื้อหาจะเปิดเผยหมด การโจมตีแบบนี้เรียกการโจมตีคนกลาง (Man-in-the-Middle) หลักการคือผู้โจมตีแทรกตัวระหว่างคุณกับเราเตอร์ Wi-Fi จริง ที่น่ากลัวกว่านั้น มือถือหลายคนตั้งเชื่อมต่ออัตโนมัติกับ Wi-Fi ที่รู้จัก ผู้โจมตีแค่ตั้งชื่อฮอตสปอตตัวเองเป็นชื่อ Wi-Fi ที่คุณเคยต่อ (เช่น ChinaNet หรือ CMCC) มือถือคุณจะเชื่อมต่อเองทันที

แล้วกันยังไง? คำแนะนำใช้ได้จริงสักข้อ:

\begin{itemize}
  \item \textbf{ยึด HTTPS ให้ขึ้นใจ}: หน้าแถบที่อยู่ขึ้นต้นด้วย https:// ตัว s พ่วงมาแปลว่า Secure ตอนเข้าเว็บไซต์ที่ต้องกรอกรหัสผ่าน ถ้าเห็น http:// (ไม่มี s) หยุดทันทีอย่าพิมพ์! เบราว์เซอร์กระแสหลักตอนนี้จะโชว์คำเตือนไม่ปลอดภัยหน้าเว็บที่ไม่ใช่ HTTPS เห็นป้ายนี้ต้องระวัง
  \item \textbf{ปิดเชื่อมต่ออัตโนมัติ}: ไม่ว่า iPhone หรือแอนดรอยด์ ไปที่การตั้งค่า Wi-Fi ปิดเข้าร่วมอัตโนมัติหรือเชื่อมต่ออัตโนมัติ ยอมกดมือทีละครั้ง ดีกว่าปล่อยให้มือถือตัดสินใจเอง
  \item \textbf{ใช้ดาต้ามือถือทำงานอ่อนไหว}: โอนเงิน ล็อกอินธนาคาร กรอกเลขบัตรประชาชน… พวกนี้ใช้ดาต้า 4G/5G ของตัวเองอย่างซื่อสัตย์ อย่าประหยัดค่าดาต้าจำนวนน้อยๆ
  \item \textbf{ลืม Wi-Fi ที่ไม่ค่อยใช้}: เดินทางกลับมาจากต่างจังหวัดหรือเที่ยวกลับมา เอา Wi-Fi โรงแรม Wi-Fi สนามบินที่เคยต่อชั่วคราวลืมทิ้งให้หมด เพื่อครั้งหน้ามือถือจะไม่เชื่อมต่อเองกับฮอตสปอตปลอมที่ตั้งชื่อเหมือนกัน
\end{itemize}

\booksection{สาม HTTPS ป้องกันคุณยังไง? — การจับมือ SSL/TLS}

เมื่อกี้พูดถึง HTTPS แล้วมันทำงานยังไงจริงๆ? ทุกครั้งที่คุณเปิดเว็บ HTTPS เบราว์เซอร์กับเซิร์ฟเวอร์จะจับมือ SSL/TLS handshake กันหนึ่งรอบ กระบวนการโดยรวมประมาณนี้:

\begin{enumerate}
  \item เบราว์เซอร์คุณบอกว่า: เฮ้ย เซิร์ฟเวอร์ ฉันอยากคุยกับนายแบบปลอดภัย นี่คือรายการความสามารถเข้ารหัสของฉัน (อัลกอริทึมเข้ารหัสที่รองรับ ฯลฯ)
  \item เซิร์ฟเวอร์ตอบกลับ: โอเค ฉันเลือกแผนเข้ารหัสชุดนี้จากรายการของนาย นี่คือใบรับรองดิจิทัลของฉัน ข้างในมีกุญแจสาธารณะ ลงนามรับรองโดยหน่วยงาน CA อะไรอะไรนั่น
  \item เบราว์เซอร์ตรวจสอบความจริงของใบรับรอง: ออกให้โดย CA ที่น่าเชื่อถือไหม? หมดอายุหรือยัง? โดเมนตรงกันไหม?
  \item ผ่านการตรวจแล้ว ทั้งสองฝ่ายใช้กุญแจสาธารณะเข้ารหัสเจรจากุญแจประจำเซสชันออกมา การสื่อสารต่อจากนี้เข้ารหัสด้วยกุญแจตัวนี้ทั้งหมด
\end{enumerate}

เหมือนคุณไปทำธุรกรรมที่ธนาคาร: คุณยื่นบัตรประชาชน (ใบรับรอง) ก่อน พนักงานตรวจยืนยันตัวตนคุณ แล้วสองคนนัดรหัสลับกัน บทสนทนาต่อจากนี้มีแค่สองคนเข้าใจ แม้มีคนแอบฟังอยู่ข้างๆ ก็ได้ยินแค่ตัวอักษรเข้ารหัสที่ระเกะระกะ

แต่จำไว้: HTTPS การันตีแค่ว่ากระบวนการส่งถูกเข้ารหัส ไม่การันตีว่าเว็บเป้าหมายเป็นคนดี เว็บฟิชชิงก็ใช้ HTTPS ได้ — ไอคอนกุญแจเล็กๆ บอกได้แค่ว่าคุณเชื่อมกับเว็บไซต์สักเว็บ ไม่ได้แปลว่าเป็นเว็บที่คุณอยากไป ต้องเช็ก URL เองว่าถูกต้อง

\booksection{สี่ ความปลอดภัย DNS: สมุดโทรศัพท์เครือข่ายของคุณปลอดภัยไหม?}

DNS (Domain Name System) เหมือนสมุดโทรศัพท์ของอินเทอร์เน็ต: คุณพิมพ์ baidu.com DNS ช่วยค้นว่ามันตรงกับ IP ไหน (เช่น 110.242.68.66) แล้วเบราว์เซอร์ถึงจะเชื่อมไปได้ ปัญหาคือการค้น DNS แบบดั้งเดิมเป็น\textbf{ข้อความล้วน}และ\textbf{ไม่มีการเข้ารหัส}! เท่ากับว่า ISP ผู้ให้ Wi-Fi หรือโหนดกลางใดๆ ก็เห็นว่าคุณเข้าเว็บไหนบ้าง

น่าเบื่อกว่านั้นคือการแอบแก้ DNS: ผู้โจมตีทำมือที่ขั้นตอน DNS ส่ง baidu.com ไปที่เว็บฟิชชิงปลอม คุณพิมพ์ URL ถูกต้อง แต่ถูกพาไปหน้าล็อกอินที่เหมือนมากแต่จริงๆ ปลอม พอกรอกบัญชีรหัสผ่านก็จบ

ดีที่ตอนนี้มีแผน DNS เข้ารหัสแล้ว:

\begin{itemize}
  \item \textbf{DNS over HTTPS (DoH)}: ซ่อนการค้นหา DNS ไว้ในดาต้า HTTPS หลอมอยู่กับดาต้าเปิดเว็บปกติ แยกไม่ออกเลยว่าอันไหนคือการค้น DNS
  \item \textbf{DNS over TLS (DoT)}: ใช้โปรโตคอลเข้ารหัส TLS คุ้มครองการค้น DNS โดยเฉพาะ พอร์ตแยก ผู้ดูแลเครือข่ายรู้ว่าคุณใช้ DoT แต่มองไม่เห็นเนื้อหา
\end{itemize}

ฉันส่วนตัวแนะนำใช้ Cloudflare 1.1.1.1 หรือ Google 8.8.8.8 เป็นเซิร์ฟเวอร์ DNS ทั้งคู่รองรับ DNS เข้ารหัส ในการตั้งค่า Private DNS ของแอนดรอยด์กรอก one.one.one.one หรือ dns.google ตรงๆ ก็จบ ง่ายมาก iPhone ยุ่งกว่าหน่อย ต้องติดตั้งโปรไฟล์หรือตั้งผ่านแอป แต่ก็คุ้มที่จะทำ

\booksection{ห้า โหมดไม่ระบุตัวตนล่องหนได้จริงไหม?}

เอาข้อสรุปก่อน:\textbf{โหมดไม่ระบุตัวตนคือม่านบ้านคุณ ไม่ใช่ผ้าคลุมล่องหน}

มันทำอะไร: ไม่ทิ้งร่องรอยในเครื่อง — ปิดหน้าต่างไม่ระบุตัวตนแล้ว เบราว์เซอร์ไม่เก็บประวัติการเข้าเว็บ Cookies บันทึกการกรอกฟอร์ม ประวัติค้นหาของคุณ ยืมคอมคนอื่นค้นข้อมูลสักอย่าง หรือซื้อของขวัญไม่อยากให้ครอบครัวเจอจากประวัติเว็บ ใช้โหมดไม่ระบุตัวตนเหมาะสม

มันทำไม่ได้อะไร:

\begin{itemize}
  \item เว็บที่คุณเข้า\textbf{ยังเห็นว่าคุณมาอยู่}, IP ของคุณ ข้อมูลอุปกรณ์ พฤติกรรมการเข้าเว็บ เว็บฝั่งนั้นเห็นหมดเลย
  \item ผู้ให้บริการเครือข่ายของคุณ\textbf{ยังเห็นว่าคุณเข้าเว็บไหนบ้าง}แค่มองไม่เห็นเนื้อหาเฉพาะ (ถ้าใช้ HTTPS)
  \item ผู้ดูแลเครือข่ายบริษัท/โรงเรียน\textbf{เฝ้าดูพฤติกรรมเข้าเน็ตของคุณได้เหมือนเดิม}ไฟล์ที่ดาวน์โหลด บุ๊กมาร์กที่เพิ่มจะ
  \item ค้างอยู่\textbf{;}ดังนั้นจำไว้: โหมดไม่ระบุตัวตนคือเครื่องลบภายในเครื่อง ไม่ใช่ชุดล่องหนบนเน็ต ถ้าจริงๆ ต้องซ่อนร่องรอยการเข้าเน็ต ต้องใช้คอมโบ VPN+Tor แต่สำหรับการใช้งานประจำวัน รู้เท่านี้พอแล้ว
\end{itemize}

หก ความเป็นส่วนตัวเบราว์เซอร์: อย่าให้เบราว์เซอร์ขายคุณ

\booksection{ความเป็นส่วนตัวของเบราว์เซอร์: อย่าให้เบราว์เซอร์ขายคุณ}

เคลียร์ข้อมูลการเข้าเว็บเป็นระยะ

\begin{itemize}
  \item \textbf{: ทุกช่วงเวลาหนึ่ง ไปตั้งค่าเคลียร์ Cookies แคช ประวัติ ไม่ใช่เพื่อกันการติดตาม — การล้าง Cookies แก้เรื่องนั้นไม่ค่อยได้ผลแล้ว — แต่เพื่อกันมือถือกระตุกขึ้นเรื่อยๆ แคชเบราว์เซอร์กินพื้นที่ได้หลาย GB}เลิกความฝันเรื่องโปรดอย่าติดตาม
  \item \textbf{: ตัวเลือกส่งคำขอ Do Not Track ในเบราว์เซอร์ แท้จริงคือขอร้องเว็บว่าอย่าติดตามคุณ แต่เว็บจะฟังหรือไม่เป็นอีกเรื่อง ส่วนใหญ่แล้วเมินทันที อย่าคาดหวังกับมัน}ใช้เบราว์เซอร์ที่คุ้มครองความเป็นส่วนตัวให้เป็น
  \item \textbf{: Firefox Focus, เบราว์เซอร์ DuckDuckGo, Brave พวกนี้บล็อกตัวติดตามกับโฆษณาเป็นค่าเริ่มต้น สะอาดกว่าการตั้งค่าเริ่มต้นของ Chrome กับ Safari มาก}เครื่องมือค้นหาก็สำคัญ
  \item \textbf{: ใช้ DuckDuckGo แทน Baidu/Google มันไม่บันทึกประวัติการค้นหาของคุณ ไม่สร้างภาพร่างผู้ใช้ พูดตรงๆ ผลการค้นหาไม่แม่นเท่า Google แต่พอสำหรับงานประจำวัน}ตรวจสิทธิ์แอป
  \item \textbf{: ไปการตั้งค่ามือถือดูว่าเบราว์เซอร์คุณได้รับสิทธิ์อะไรไว้ มันจำเป็นจริงๆ ที่ต้องเข้าถึงตำแหน่ง ไมค์ สมุดโทรศัพท์ของคุณเหรอ? โอกาสสูงมากที่ไม่จำเป็น}เจ็ด เช็กลิสต์ความปลอดภัยที่ใช้ได้จริง: คุ้มครองตัวเองตั้งแต่วันนี้
\end{itemize}

\booksection{เช็คลิสต์ความปลอดภัยที่ใช้งานได้จริง: เริ่มปกป้องตัวเองตั้งแต่วันนี้}

ปิดเชื่อมต่อ Wi-Fi อัตโนมัติ

\begin{enumerate}
  \item \textbf{: iPhone อยู่ใน ตั้งค่า-Wi-Fi แอนดรอยด์อยู่ในการตั้งค่า Wi-Fi}เปลี่ยน DNS แบบเข้ารหัส
  \item \textbf{: แอนดรอยด์ตั้ง Private DNS เป็น dns.google ตรงๆ iPhone ใช้แอปหรือโปรไฟล์ตั้ง DoH ของ 1.1.1.1}เปิดการยืนยันสองปัจจัย (2FA)
  \item \textbf{: WeChat Alipay Apple ID บัญชี Google อะไรเปิด 2FA ได้ก็เปิด รหัสยืนยันทาง SMS ถือเป็น 2FA ระดับต่ำสุดแล้ว ใช้ TOTP (เช่น Google Authenticator, Microsoft Authenticator) ยิ่งดี}ไล่ตรวจสิทธิ์แอปรอบหนึ่ง
  \item \textbf{: ไล่ดูสิทธิ์ของแอปทีละตัว โดยเฉพาะห้าอย่างนี้: ตำแหน่ง สมุดโทรศัพท์ SMS ไมค์ กล้อง คุณจะเจอแอปไฟฉายบางตัวขอสมุดโทรศัพท์คุณด้วยซ้ำ — ถอนทิ้งเลย}เคลียร์ข้อมูลเบราว์เซอร์เป็นระยะ
  \item \textbf{: ตั้งการเตือนปฏิทินไว้ เดือนละครั้ง ล้างแคช Cookies}อย่าใช้สายข้อมูลของชาร์จสาธารณะ
  \item \textbf{: พอร์ต USB ชาร์จได้แต่ก็ถ่ายโอนข้อมูลได้ ชาร์จประสงค์ร้าย (เรียก Juice Jacking) แอบเอาข้อมูลคุณตอนชาร์จได้ ออกเรือนพกหัวชาร์จกับสายของตัวเอง}ทำเลยตอนนี้: ทำตรวจสุขภาพความปลอดภัยให้มือถือคุณ
\end{enumerate}

\begin{tcbpractice}{อัปเกรด DNS}
\begin{enumerate}
  \item \textbf{: แอนดรอยด์: ตั้งค่า > การเชื่อมต่อและการแชร์ > Private DNS > กรอก one.one.one.one หรือ dns.google iPhone: ใช้ Safari ค้น 1.1.1.1 DNS profile ติดตั้งโปรไฟล์ หรือดาวน์โหลดแอป 1.1.1.1 ของ Cloudflare}ตรวจสอบ Wi-Fi
  \item \textbf{: เปิดรายการ Wi-Fi เอาฮอตสปอตที่ไม่รู้จัก ไม่ค่อยใช้ ชื่อแปลกๆ ลืมทิ้งหมด}กวาดสิทธิ์ใหญ่
  \item \textbf{: เข้า ตั้งค่า > ความเป็นส่วนตัว/การจัดการสิทธิ์ ไล่ตรวจสิทธิ์ทีละแอป จดแอปที่ขอสิทธิ์ไม่สมเหตุสมผล (เช่นเครื่องคิดเลขขอสมุดโทรศัพท์) พิจารณาถอนหรือหาตัวแทน}ตรวจความปลอดภัยบัญชี
  \item \textbf{: เปิดหน้าความปลอดภัยของ Apple ID หรือบัญชี Google ดูว่ามีอุปกรณ์แปลกปลอมล็อกอินไหม ยืนยันว่าเปิด 2FA แล้ว ตรวจดูว่าผู้ติดต่อกู้คืนน่าเชื่อถือไหม}ขั้นต่อไป: AI บนเครื่องเก็บการประมวลผลไว้ในมือถือ ทั้งเร็ว ปลอดภัยต่อความเป็นส่วนตัว และใช้แบบออฟไลน์ได้ ส่วน AI บนคลาวด์ต้องต่อเน็ต แต่รับมือกับงานซับซ้อนกว่าได้ (เจาะลึกในบทที่ 22)
\end{enumerate}

\end{tcbpractice}
